\section{Limitações e agenda de pesquisa}
O corpus é pequeno, dirigido por problemas e desigual no tempo. As lacunas de várias décadas do século XX impedem estabelecer predomínio de formas, taxas de substituição ou uma cronologia nacional de difusão. A coleta não autenticou a série de 22 transcrições enviada inicialmente e não identificou a primeira ocorrência absoluta de \forma{kkkk}. Esses limites são resultados da auditoria, não espaços a preencher com frases plausíveis.

O acesso digital introduz filtros: preservação seletiva, qualidade de digitalização, OCR, disponibilidade de catálogo, direitos e barreiras técnicas. Uma ausência num acervo pode refletir qualquer desses fatores. A data de uma transcrição eletrônica não é a data da criação literária; a data de publicação de uma dissertação não é necessariamente a data de uma mensagem. O estudo registra essas diferenças, mas não resolve todos os problemas de transmissão editorial dos textos consultados.

A agenda de ampliação tem quatro frentes. Primeiro, construir uma moldura anual de edições e números de periódicos, com países, gêneros e extensões. Segundo, buscar registros digitais contemporâneos anteriores a 2007, validando procedência e data, sem confundir memórias recentes com logs antigos. Terceiro, estudar canais intermediários como teatro, quadrinhos, cartas e livros escolares, capazes de testar continuidade indireta. Quarto, investigar a pragmática contemporânea com dados brasileiros e variáveis sociais observadas, incluindo comparações entre dispositivos e comunidades de prática.

A aplicação do plano estatístico dependerá do objetivo: existência, detecção de formas raras ou estimação de frequência. As quotas de três livros e três periódicos podem permanecer como seleção de exemplos para leitura; não devem ser apresentadas como amostra representativa. Um estudo probabilístico requer sorteio, denominadores e análise compatíveis. Um estudo filológico pode ter grande valor com poucos testemunhos, desde que não extrapole sua abrangência.

\section{Conclusão}
Os testemunhos recuperados confirmam diversidade de formas de escrever o riso, incluindo \forma{quiá}, \forma{cá}, \forma{ah} e \forma{quá}, além de cadeias digitais de \forma{k}. O registro de \textit{O Tico-Tico}, de 1935, fornece uma localização verificável de \forma{quá}; a dissertação de Young oferece documentação de \forma{kkkk} em pesquisa situada em 2007.2. Esses achados sustentam atestações específicas, sem determinar a origem absoluta da convenção digital.

A contribuição conceitual é rejeitar duas inferências simétricas. A semelhança entre representações não demonstra uma genealogia; o caráter onomatopaico e as lacunas documentais não demonstram independência. Onomatopeias podem integrar a história das línguas, e mudanças gráficas podem ter consequências linguísticas. A questão é empírica: quais canais, documentos e regularidades sustentam uma relação proposta?

O dimensionamento estatístico mostra por que poucas referências por período servem melhor à exemplificação que à prova de ausência ou predomínio. A documentação pública, as tabelas de cobertura e os scripts permitem ampliar e corrigir a pesquisa. A interpretação cientificamente defensável permanece aberta à descoberta de novos registros, inclusive anteriores aos aqui localizados.

\section*{Disponibilidade de dados e materiais}
Manuscrito, bibliografia, inventários, códigos de evidência e scripts: \repo. Os downloads permitidos para distribuição acompanham o repositório; arquivos com condições de republicação não verificadas permanecem no acervo local de trabalho, com links públicos de aquisição. O arquivo \texttt{README.md} descreve a reprodução e as lacunas. A versão consultada deve ser identificada por commit ou release.

\section*{Declarações}
\textbf{Autoria:} Murillo Edson de Carvalho Souza, único autor. \textbf{Financiamento e conflitos de interesse:} não foram informados no material fornecido; nenhuma declaração de inexistência é presumida. \textbf{Assistência automatizada:} utilizada na elaboração e documentação, conforme descrito na discussão. \textbf{Situação editorial:} manuscrito de pesquisa exploratória, sem indicação de avaliação por pares concluída.
